\section{Polynomial-Block Contraction Criterion}
\label{sec:block-contraction}

A spectral-radius bound alone is not sufficient to prove a polynomial hitting
time for a general nonnormal substochastic transition matrix.  For example, a
long nilpotent shift may have spectral radius zero while requiring a number of
steps proportional to the number of transient states.  A stronger and directly
probabilistic certificate is therefore used.

Let $Q_F$ be the transition matrix restricted to states that have not yet
reached the target set.  For satisfiable formulas the target is the set of
zero-energy assignments.  For first-passage analysis of an unsatisfiable
formula, the target may instead be the set of global minima.

For an integer block length $L$, define
\begin{equation}
s_L(F)
=
\left\|Q_F^L\mathbf 1\right\|_\infty
\end{equation}
and
\begin{equation}
\varepsilon_L(F)=1-s_L(F).
\end{equation}
Here $s_L(F)$ is the largest, over all starting states, probability of avoiding
the target for the next $L$ updates.

\begin{theorem}[Geometric block bound]
If
\begin{equation}
\varepsilon_L(F)>0,
\end{equation}
then the expected first-passage time satisfies
\begin{equation}
\boxed{
\max_x\mathbb E_x[T]
\le
\frac{L}{\varepsilon_L(F)}.
}
\end{equation}
\end{theorem}

\begin{proof}
Conditioned on any history that has not yet reached the target, the probability
of reaching the target during the next $L$ updates is at least
$\varepsilon_L(F)$.  Therefore the number of $L$-step blocks required for
success is stochastically dominated by a geometric random variable of mean
$1/\varepsilon_L(F)$.
\end{proof}

Consequently, a uniform theorem of the form
\begin{equation}
L(n,m)=\operatorname{poly}(n+m),
\qquad
\varepsilon_{L(n,m)}(F)
\ge
\frac{1}{\operatorname{poly}(n+m)}
\end{equation}
for every 3-CNF formula would imply polynomial expected first-passage time.

An evolutionary adversary was used to minimize the block absorption margin at
the polynomial block length $L=n^2$.

\begin{table}[t]
\centering
\caption{Adversarial $n^2$-block contraction search.}
\begin{tabular}{rrrrrr}
\toprule
$n$ & $L$ & Max survival & $\varepsilon_L$ & $L/\varepsilon_L$ &
$\max H_F$ \\
\midrule
8 & 64 & 0.826229 & 0.173771 & 368.30 & 275.00 \\
10 & 100 & 0.899455 & 0.100545 & 994.58 & 738.99 \\
12 & 144 & 0.752752 & 0.247248 & 582.41 & 357.17 \\
\bottomrule
\end{tabular}
\end{table}

The search did not produce a vanishing block margin at these small sizes.  This
is finite-instance evidence only and does not establish a uniform lower bound
on $\varepsilon_L(F)$.
